\section{A Minimal, Deterministic Kernel}
\label{sec:kernel}

The conservation law (\cref{thm:conservation}) describes \emph{what} the
kernel computes. This section describes \emph{how}, and why the
how is fast precisely \emph{because} the what is small.

\subsection{Non-expressiveness as a design goal}

The kernel's sole responsibility is to enforce the mechanical
invariants: that a transaction conserves value, and that the ledger is
the deterministic fold of the log. It encodes no posting rules, no
account hierarchy semantics, no interpretation policy, and exposes no
mutable configuration. This is a deliberate inversion of the usual
priority. Most systems make the core expressive so that business needs
can be met ``in the engine''; Ratio makes the core \emph{non-expressive}
so that the engine is small enough to prove and cheap enough to run, and
pushes expressiveness outward into the control plane.

Non-expressiveness pays off twice. For \emph{assurance}, the trusted
computing base is a handful of definitions and the theorems over them ---
the surface a reviewer (or a proof assistant) must trust is tiny. For
\emph{performance}, a kernel that does not branch on business semantics
does not branch much at all: it streams fixed-width records and adds
integers.

\subsection{Data layout and the cost model}

A transaction's postings are stored as cache-aligned, fixed-width binary
records: an integer dimension index and an \texttt{i64} minor-unit
amount, never a tagged union of value types and never floating point.
The ledger fold (\cref{def:ledger}) is then a branch-minimal loop ---
load record, accumulate into a per-dimension integer, advance --- with no
data-dependent control flow on the hot path. Its arithmetic is exact by
construction (\cref{sec:proofs}), so there is no reconciliation pass and
no tolerance epsilon to tune.

The consequence is a cost model dominated by \emph{memory bandwidth}
rather than control-flow complexity. Throughput scales with how fast the
machine can stream the log past the ALU, not with how many business
cases the engine must dispatch on. \Cref{fig:throughput} contrasts the
asymptotics: an expressive engine pays a per-transaction dispatch and
allocation cost that grows with rule-set complexity, while the
conservation kernel's per-transaction cost is a small constant set by the
record width. \Cref{fig:determinism} makes the determinism claim
concrete: because the inputs are immutable and the fold is fixed and
integer-only, independent replays of the same log produce bit-identical
ledger states --- the variance across runs is exactly zero, by
construction, where a float-based or order-sensitive engine accumulates
drift.

\subsection{Determinism and replay}

Three properties combine to make any ledger state exactly reconstructible:
the log is append-only, the fold is a fixed monoidal aggregation, and all
arithmetic is over $\mathbb{Z}$. There is no source of nondeterminism on
the path from log to state --- no floating-point rounding, no hash-map
iteration order leaking into a sum, no wall-clock dependence. Replaying a
prefix of the log reproduces the historical state it implies, which is
what makes point-in-time audit and ``what did we believe on date $T$''
queries first-class rather than forensic.

\Cref{fig:architecture} places the kernel in the running system: the
gRPC \texttt{Ledger} service (an AIP-idiomatic~\cite{aip} contract --- \texttt{GetTransaction},
\texttt{ListTransactions}, \texttt{CreateTransaction}) is the only door to
the kernel, and \texttt{CreateTransaction} rejects any transaction that
fails conservation with \texttt{FAILED\_PRECONDITION}. The kernel's
invariant is thus enforced at the system boundary by the same predicate
the proofs reason about.

% System architecture. The diagram itself is authored in PlantUML
% (figures/fig_architecture.puml) and rendered to vector PDF by
% rules_puml (PlantUML + Apache Batik/FOP, no external toolchain);
% this wrapper embeds the rendered PDF.
\begin{figure}[t]
\centering
\includegraphics[width=\linewidth]{figures/fig_architecture_diagram.pdf}
\caption{Ratio in the running system. Natural-language intent is turned
into a candidate specification by an LLM and compiled, deterministically,
into a content-addressed control-plane artifact (top, untrusted). The
only path to ledger state is the gRPC \texttt{Ledger} contract, whose
\texttt{CreateTransaction} admits a transaction \emph{iff} it conserves
value --- the kernel's sole invariant, proved in Lean and emitted to the
Rust that runs. Learned and declarative components sit outside the
trusted core (green) and cannot violate conservation.}
\label{fig:architecture}
\end{figure}

